% Options for packages loaded elsewhere
\PassOptionsToPackage{unicode}{hyperref}
\PassOptionsToPackage{hyphens}{url}
\documentclass[
  slovene,
]{article}
\usepackage{xcolor}
\usepackage{amsmath,amssymb}
\setcounter{secnumdepth}{-\maxdimen} % remove section numbering
\usepackage{iftex}
\ifPDFTeX
  \usepackage[T1]{fontenc}
  \usepackage[utf8]{inputenc}
  \usepackage{textcomp} % provide euro and other symbols
\else % if luatex or xetex
  \usepackage{unicode-math} % this also loads fontspec
  \defaultfontfeatures{Scale=MatchLowercase}
  \defaultfontfeatures[\rmfamily]{Ligatures=TeX,Scale=1}
\fi
\usepackage{lmodern}
\ifPDFTeX\else
  % xetex/luatex font selection
\fi
% Use upquote if available, for straight quotes in verbatim environments
\IfFileExists{upquote.sty}{\usepackage{upquote}}{}
\IfFileExists{microtype.sty}{% use microtype if available
  \usepackage[]{microtype}
  \UseMicrotypeSet[protrusion]{basicmath} % disable protrusion for tt fonts
}{}
\makeatletter
\@ifundefined{KOMAClassName}{% if non-KOMA class
  \IfFileExists{parskip.sty}{%
    \usepackage{parskip}
  }{% else
    \setlength{\parindent}{0pt}
    \setlength{\parskip}{6pt plus 2pt minus 1pt}}
}{% if KOMA class
  \KOMAoptions{parskip=half}}
\makeatother
\ifLuaTeX
\usepackage[bidi=basic,shorthands=off]{babel}
\else
\usepackage[bidi=default,shorthands=off]{babel}
\fi
\ifLuaTeX
  \usepackage{selnolig} % disable illegal ligatures
\fi
\setlength{\emergencystretch}{3em} % prevent overfull lines
\providecommand{\tightlist}{%
  \setlength{\itemsep}{0pt}\setlength{\parskip}{0pt}}
\usepackage{bookmark}
\IfFileExists{xurl.sty}{\usepackage{xurl}}{} % add URL line breaks if available
\urlstyle{same}
% fallback for those not using the hyperref driver hyperxmp:
\makeatletter
\@ifundefined{xmpquote}{\newcommand{\xmpquote}[1]{#1}}{}
\makeatother
\hypersetup{
  pdftitle={Fourierove vrste{,} Baselski problem in Gaussov integral},
  pdflang={sl-SI},
  hidelinks,
  pdfcreator={LaTeX via pandoc}}

\title{Fourierove vrste, Baselski problem in Gaussov integral}
\author{}
\date{}

\begin{document}
\maketitle

\section{Ortogonalnost trigonometrijskih
funkcij}\label{ortogonalnost-trigonometrijskih-funkcij}

Naj bosta \(m, n \in \mathbb{N}\). Za vse take \(m\) in \(n\) veljajo
zveze

\[\int_{-\pi}^{\pi} \sin(mx) \sin(nx) \, dx = \begin{cases} \pi, & m = n, \\ 0, & m \neq n, \end{cases}\]

\[\int_{-\pi}^{\pi} \cos(mx) \cos(nx) \, dx = \begin{cases} \pi, & m = n, \\ 0, & m \neq n, \end{cases}\]

\[\int_{-\pi}^{\pi} \sin(mx) \cos(nx) \, dx = 0.\]

Prvo zvezo dokažemo z adicijskim izrekom. Ker je
\(\sin \alpha \sin \beta = \frac{1}{2} \bigl( \cos(\alpha - \beta) - \cos(\alpha + \beta) \bigr)\),
je za \(m \neq n\)

\[\begin{aligned}
\int_{-\pi}^{\pi} \sin(mx) \sin(nx) \, dx
&= \frac{1}{2} \int_{-\pi}^{\pi} \bigl( \cos((m - n)x) - \cos((m + n)x) \bigr) \, dx \\
&= \frac{1}{2} \left[ \frac{\sin((m - n)x)}{m - n} - \frac{\sin((m + n)x)}{m + n} \right]_{-\pi}^{\pi} \\
&= 0,
\end{aligned}\]

saj je \(\sin(k\pi) = 0\) za vsak \(k \in \mathbb{Z}\). Za \(m = n\) pa
je

\[\int_{-\pi}^{\pi} \sin^2(nx) \, dx = \frac{1}{2} \int_{-\pi}^{\pi} \bigl( 1 - \cos(2nx) \bigr) \, dx = \frac{1}{2} \cdot 2\pi = \pi.\]

Drugo in tretjo zvezo dokažemo enako, z identitetama

\[\cos \alpha \cos \beta = \frac{1}{2} \bigl( \cos(\alpha - \beta) + \cos(\alpha + \beta) \bigr),\]

\[\sin \alpha \cos \beta = \frac{1}{2} \bigl( \sin(\alpha - \beta) + \sin(\alpha + \beta) \bigr).\]

\section{Fourierova vrsta}\label{fourierova-vrsta}

Naj bo \(f \colon [-\pi, \pi] \to \mathbb{R}\) integrabilna funkcija.
Njena \emph{Fourierova vrsta} je

\[f(x) \sim \frac{a_0}{2} + \sum_{n=1}^{\infty} \bigl( a_n \cos(nx) + b_n \sin(nx) \bigr),\]

kjer so koeficienti enaki

\[a_n = \frac{1}{\pi} \int_{-\pi}^{\pi} f(x) \cos(nx) \, dx, \qquad b_n = \frac{1}{\pi} \int_{-\pi}^{\pi} f(x) \sin(nx) \, dx.\]

Formuli za koeficienta dobimo iz ortogonalnosti: vrsto pomnožimo s
\(\cos(mx)\) oziroma s \(\sin(mx)\) in jo integriramo po členih. Od vseh
členov ostane samo tisti z indeksom \(n = m\).

\subsection{Pravokotni signal}\label{pravokotni-signal}

Za zgled vzemimo liho funkcijo

\[f(x) = \operatorname{sgn} x = \begin{cases} -1, & -\pi < x < 0, \\ 0, & x = 0, \\ 1, & 0 < x < \pi. \end{cases}\]

Ker je \(f\) liha, je \(f(x) \cos(nx)\) liha funkcija, zato je
\(a_n = 0\) za vse \(n \geq 0\). Funkcija \(f(x) \sin(nx)\) je soda,
zato je

\[\begin{aligned}
b_n &= \frac{2}{\pi} \int_{0}^{\pi} \sin(nx) \, dx
= \frac{2}{\pi} \left[ -\frac{\cos(nx)}{n} \right]_{0}^{\pi}
= \frac{2}{n\pi} \bigl( 1 - \cos(n\pi) \bigr)
= \frac{2}{n\pi} \bigl( 1 - (-1)^n \bigr) \\
&= \begin{cases} \dfrac{4}{n\pi}, & n \text{ lih}, \\ 0, & n \text{ sod}. \end{cases}
\end{aligned}\]

Fourierova vrsta pravokotnega signala je torej

\[\begin{aligned}
\operatorname{sgn} x &= \frac{4}{\pi} \sum_{k=0}^{\infty} \frac{\sin\bigl((2k + 1)x\bigr)}{2k + 1} \\
&= \frac{4}{\pi} \left( \sin x + \frac{\sin(3x)}{3} + \frac{\sin(5x)}{5} + \cdots \right), \qquad -\pi < x < \pi.
\end{aligned}\]

Vrsta konvergira k \(f(x)\) v vsaki točki, v kateri je \(f\) zvezna, v
točki \(x = 0\) pa k povprečju leve in desne limite,
\(\frac{1}{2} \bigl( f(0^-) + f(0^+) \bigr) = 0\).

\subsection{Leibnizova vrsta}\label{leibnizova-vrsta}

Vstavimo \(x = \frac{\pi}{2}\). Ker je
\(\sin\bigl((2k + 1)\frac{\pi}{2}\bigr) = (-1)^k\), dobimo

\[1 = \frac{4}{\pi} \sum_{k=0}^{\infty} \frac{(-1)^k}{2k + 1}, \qquad \text{torej} \qquad \frac{\pi}{4} = 1 - \frac{1}{3} + \frac{1}{5} - \frac{1}{7} + \cdots = \sum_{k=0}^{\infty} \frac{(-1)^k}{2k + 1}.\]

\section{Parsevalova enakost in Baselski
problem}\label{parsevalova-enakost-in-baselski-problem}

Za kvadratno integrabilno funkcijo \(f\) velja \emph{Parsevalova
enakost}

\[\frac{1}{\pi} \int_{-\pi}^{\pi} \lvert f(x) \rvert^2 \, dx = \frac{a_0^2}{2} + \sum_{n=1}^{\infty} \bigl( a_n^2 + b_n^2 \bigr).\]

Za pravokotni signal je \(\lvert f(x) \rvert^2 = 1\) povsod razen v
točki \(0\), zato je leva stran enaka \(\frac{1}{\pi} \cdot 2\pi = 2\).
Na desni ostanejo samo lihi koeficienti
\(b_{2k+1} = \frac{4}{(2k + 1)\pi}\):

\[2 = \sum_{k=0}^{\infty} \frac{16}{\pi^2 (2k + 1)^2}, \qquad \text{torej} \qquad \sum_{k=0}^{\infty} \frac{1}{(2k + 1)^2} = \frac{\pi^2}{8}.\]

Naj bo \(S = \sum_{n=1}^{\infty} \frac{1}{n^2}\). Vsoto razdelimo na
sode in lihe člene:

\[\begin{aligned}
S &= \sum_{k=1}^{\infty} \frac{1}{(2k)^2} + \sum_{k=0}^{\infty} \frac{1}{(2k + 1)^2} \\
&= \frac{1}{4} \sum_{k=1}^{\infty} \frac{1}{k^2} + \frac{\pi^2}{8} \\
&= \frac{S}{4} + \frac{\pi^2}{8}.
\end{aligned}\]

Od tod je \(\frac{3}{4} S = \frac{\pi^2}{8}\) in

\[\boxed{\sum_{n=1}^{\infty} \frac{1}{n^2} = \frac{\pi^2}{6}.}\]

Na enak način dobimo še vsoto alternirajoče vrste:

\[\sum_{n=1}^{\infty} \frac{(-1)^{n+1}}{n^2} = \sum_{k=0}^{\infty} \frac{1}{(2k + 1)^2} - \sum_{k=1}^{\infty} \frac{1}{(2k)^2} = \frac{\pi^2}{8} - \frac{1}{4} \cdot \frac{\pi^2}{6} = \frac{\pi^2}{12}.\]

\section{Gaussov integral}\label{gaussov-integral}

Izračunajmo

\[I = \int_{-\infty}^{\infty} e^{-x^2} \, dx.\]

Nedoločenega integrala funkcije \(e^{-x^2}\) ne moremo zapisati z
elementarnimi funkcijami, zato si pomagamo z dvojnim integralom. Po
Fubinijevem izreku je

\[I^2 = \left( \int_{-\infty}^{\infty} e^{-x^2} \, dx \right) \left( \int_{-\infty}^{\infty} e^{-y^2} \, dy \right) = \iint_{\mathbb{R}^2} e^{-(x^2 + y^2)} \, dx \, dy.\]

Preidemo na polarne koordinate \(x = r \cos \varphi\),
\(y = r \sin \varphi\). Jacobijeva determinanta te preslikave je

\[\det \frac{\partial(x, y)}{\partial(r, \varphi)} = \begin{vmatrix} \cos \varphi & -r \sin \varphi \\ \sin \varphi & r \cos \varphi \end{vmatrix} = r \cos^2 \varphi + r \sin^2 \varphi = r,\]

zato je

\[\begin{aligned}
I^2 &= \int_{0}^{2\pi} \int_{0}^{\infty} e^{-r^2} \, r \, dr \, d\varphi
= 2\pi \int_{0}^{\infty} r e^{-r^2} \, dr \\
&= 2\pi \left[ -\frac{1}{2} e^{-r^2} \right]_{0}^{\infty}
= 2\pi \cdot \frac{1}{2}
= \pi.
\end{aligned}\]

Ker je \(I > 0\), je

\[\int_{-\infty}^{\infty} e^{-x^2} \, dx = \sqrt{\pi}.\]

\subsection{Normalna porazdelitev}\label{normalna-porazdelitev}

Z uvedbo nove spremenljivke \(x = \frac{t - \mu}{\sigma \sqrt{2}}\),
kjer je \(\sigma > 0\), dobimo

\[\int_{-\infty}^{\infty} \frac{1}{\sigma \sqrt{2\pi}} \exp\left( -\frac{(t - \mu)^2}{2\sigma^2} \right) dt = \frac{1}{\sigma \sqrt{2\pi}} \cdot \sigma \sqrt{2} \cdot \sqrt{\pi} = 1,\]

zato je funkcija pod integralom gostota verjetnosti. To je gostota
normalne porazdelitve \(N(\mu, \sigma^2)\).

\subsection{Funkcija gama}\label{funkcija-gama}

Funkcijo gama za \(s > 0\) definiramo z integralom

\[\Gamma(s) = \int_{0}^{\infty} t^{s - 1} e^{-t} \, dt.\]

Z integracijo po delih dobimo \(\Gamma(s + 1) = s \, \Gamma(s)\), iz
\(\Gamma(1) = 1\) pa sledi \(\Gamma(n + 1) = n!\) za vsak
\(n \in \mathbb{N}\). Z uvedbo nove spremenljivke \(t = u^2\),
\(dt = 2u \, du\) dobimo

\[\Gamma\left( \frac{1}{2} \right) = \int_{0}^{\infty} t^{-1/2} e^{-t} \, dt = \int_{0}^{\infty} \frac{1}{u} e^{-u^2} \, 2u \, du = 2 \int_{0}^{\infty} e^{-u^2} \, du = \sqrt{\pi},\]

od tod pa na primer

\[\Gamma\left( \frac{5}{2} \right) = \frac{3}{2} \cdot \frac{1}{2} \cdot \Gamma\left( \frac{1}{2} \right) = \frac{3\sqrt{\pi}}{4}.\]

\section{Vandermondova determinanta}\label{vandermondova-determinanta}

Naj bodo \(x_1, x_2, \ldots, x_n\) realna števila. Vandermondova matrika
je

\[V = \begin{pmatrix}
1 & x_1 & x_1^2 & \cdots & x_1^{n-1} \\
1 & x_2 & x_2^2 & \cdots & x_2^{n-1} \\
\vdots & \vdots & \vdots & \ddots & \vdots \\
1 & x_n & x_n^2 & \cdots & x_n^{n-1}
\end{pmatrix},\]

njena determinanta pa je

\[\det V = \prod_{1 \leq i < j \leq n} (x_j - x_i).\]

Za \(n = 3\) formulo preverimo neposredno. Od drugega in tretjega
stolpca odštejemo \(x_1\)-kratnik prejšnjega stolpca, začnemo pri
zadnjem:

\[\begin{aligned}
\begin{vmatrix} 1 & x_1 & x_1^2 \\ 1 & x_2 & x_2^2 \\ 1 & x_3 & x_3^2 \end{vmatrix}
&= \begin{vmatrix} 1 & 0 & 0 \\ 1 & x_2 - x_1 & x_2(x_2 - x_1) \\ 1 & x_3 - x_1 & x_3(x_3 - x_1) \end{vmatrix} \\
&= (x_2 - x_1)(x_3 - x_1) \begin{vmatrix} 1 & x_2 \\ 1 & x_3 \end{vmatrix} \\
&= (x_2 - x_1)(x_3 - x_1)(x_3 - x_2).
\end{aligned}\]

Determinanta je torej različna od nič natanko takrat, ko so števila
\(x_1, \ldots, x_n\) paroma različna. V tem primeru ima sistem

\[\begin{aligned}
c_0 + c_1 x_1 + c_2 x_1^2 + \cdots + c_{n-1} x_1^{n-1} &= y_1, \\
c_0 + c_1 x_2 + c_2 x_2^2 + \cdots + c_{n-1} x_2^{n-1} &= y_2, \\
&\;\;\vdots \\
c_0 + c_1 x_n + c_2 x_n^2 + \cdots + c_{n-1} x_n^{n-1} &= y_n
\end{aligned}\]

natanko eno rešitev: skozi \(n\) točk \((x_i, y_i)\) z različnimi
absciscami poteka natanko en polinom stopnje največ \(n - 1\). Zapišemo
ga lahko v Lagrangeevi obliki

\[p(x) = \sum_{i=1}^{n} y_i \prod_{\substack{j = 1 \\ j \neq i}}^{n} \frac{x - x_j}{x_i - x_j}.\]

\end{document}
